\chapter{Θεωρητικό Μέρος}

\section{Κυβερνοασφάλεια στα Δίκτυα \en{5G} και \en{6G}}

\subsection{Εισαγωγή}

Η ταχεία εξέλιξη των κινητών δικτύων προς το \en{5G} και το επερχόμενο \en{6G} έχει εισαγάγει πρωτοφανείς δυνατότητες όσον αφορά τη συνδεσιμότητα, την ταχύτητα και τη νοημοσύνη. Ωστόσο, αυτός ο μετασχηματισμός έχει αυξήσει σημαντικά την πολυπλοκότητα των υποδομών δικτύου, οδηγώντας σε νέες προκλήσεις κυβερνοασφάλειας. Σε αντίθεση με τις προηγούμενες γενιές, τα δίκτυα επόμενης γενιάς βασίζονται σε μεγάλο βαθμό σε \en{software-defined} στοιχεία, εικονικοποίηση και κατανεμημένες αρχιτεκτονικές, τα οποία διευρύνουν την επιφάνεια επίθεσης.

Η κυβερνοασφάλεια στα δίκτυα \en{5G} και \en{6G} δεν αποτελεί απλώς επέκταση της παραδοσιακής ασφάλειας δικτύων, αλλά απαιτεί μια θεμελιωδώς νέα προσέγγιση που ενσωματώνει παρακολούθηση σε πραγματικό χρόνο, ευφυή ανίχνευση απειλών και αυτοματοποιημένους μηχανισμούς απόκρισης \cite{ferrag2019security}. Το κεφάλαιο αυτό εξετάζει τους κύριους τύπους κυβερνοαπειλών, τους παράγοντες που συμβάλλουν στη διεύρυνση της επιφάνειας επίθεσης και την ανάγκη για προηγμένες λύσεις ασφάλειας.

\subsection{Τύποι Κυβερνοαπειλών}

\subsubsection{Επιθέσεις \en{Distributed Denial of Service (DDoS)}}

Οι επιθέσεις \en{Distributed Denial of Service (DDoS)} παραμένουν μία από τις πιο διαδεδομένες απειλές στα σύγχρονα δίκτυα. Στόχος τους είναι η υπερφόρτωση των πόρων του δικτύου μέσω τεράστιου όγκου κίνησης, καθιστώντας τις υπηρεσίες μη διαθέσιμες στους νόμιμους χρήστες. Στα περιβάλλοντα \en{5G}, το υψηλό εύρος ζώνης και η μαζική συνδεσιμότητα μπορούν να ενισχύσουν τον αντίκτυπο τέτοιων επιθέσεων \cite{zhang2019ddos}.

Η ενσωμάτωση συσκευών \en{IoT} επιδεινώνει περαιτέρω το πρόβλημα, καθώς παραβιασμένες συσκευές μπορούν να σχηματίσουν μεγάλα \en{botnets}. Στο \en{6G}, η κλίμακα αυτών των επιθέσεων αναμένεται να αυξηθεί σημαντικά.

\subsubsection{Επιθέσεις Εισβολής και \en{Malware}}

Οι επιθέσεις εισβολής και \en{malware} στοχεύουν την υποδομή δικτύου και τις συσκευές χρηστών για μη εξουσιοδοτημένη πρόσβαση ή κλοπή δεδομένων. Στα δίκτυα \en{5G}, η χρήση εικονικοποίησης και \en{cloud-native} τεχνολογιών εισάγει νέες ευπάθειες, καθώς οι επιτιθέμενοι μπορούν να εκμεταλλευτούν σφάλματα λογισμικού ή λανθασμένες ρυθμίσεις σε \en{virtual network functions (VNFs)} \cite{ahmed2020survey}.

Ιδιαίτερη ανησυχία προκαλούν τα \en{advanced persistent threats (APTs)}, τα οποία παραμένουν σε ένα σύστημα για μεγάλο χρονικό διάστημα καθώς η έγκαιρη ανίχνευσή τους είναι δύσκολη.

\subsubsection{Επιθέσεις Σηματοδοσίας}

Οι επιθέσεις σηματοδοσίας εκμεταλλεύονται ευπάθειες στο επίπεδο ελέγχου των δικτύων. Δημιουργούν υπερβολικό όγκο σημάτων, υπερφορτώνοντας στοιχεία όπως το \en{Mobility Management Entity (MME)} ή το \en{Access and Mobility Management Function (AMF)} \cite{li2018signaling}.

Στα δίκτυα \en{5G}, η αρχιτεκτονική βασισμένη σε υπηρεσίες αυξάνει τις πιθανότητες εκμετάλλευσης τέτοιων αδυναμιών.

\subsubsection{Παραβιάσεις και Διαρροές Δεδομένων}

Οι παραβιάσεις και διαρροές δεδομένων αποτελούν σοβαρή απειλή, καθώς τεράστιοι όγκοι ευαίσθητων πληροφοριών μεταδίδονται και επεξεργάζονται. Η μη εξουσιοδοτημένη πρόσβαση μπορεί να έχει σοβαρές συνέπειες \cite{roman2018mobile}.

Η κατανεμημένη φύση των δικτύων \en{5G}/\en{6G}, σε συνδυασμό με \en{edge computing} και υποδομές \en{cloud}, αυξάνει τα πιθανά σημεία επίθεσης.

\subsection{Διευρυμένη Επιφάνεια Επίθεσης}

\subsubsection{Εξάπλωση \en{IoT}}

Η εξάπλωση του \en{Internet of Things (IoT)} αποτελεί βασικό παράγοντα αύξησης της επιφάνειας επίθεσης. Οι συσκευές αυτές συχνά στερούνται ισχυρών μηχανισμών ασφάλειας \cite{atzori2010internet}.

Επιπλέον, ο μεγάλος αριθμός τους αυξάνει την πιθανότητα εκμετάλλευσης ευπαθειών και δημιουργίας \en{botnets}.

\subsubsection{Εικονικοποιημένη Υποδομή}

Η υιοθέτηση τεχνολογιών όπως \en{Network Function Virtualization (NFV)} και \en{Software-Defined Networking (SDN)} εισάγει νέες προκλήσεις ασφάλειας \cite{kreutz2015software}.

Ευπάθειες σε \en{hypervisors}, εικονικές μηχανές ή πλατφόρμες ορχήστρωσης μπορούν να επηρεάσουν πολλαπλές λειτουργίες ταυτόχρονα.

\subsection{Ανάγκη για Ανίχνευση και Απόκριση σε Πραγματικό Χρόνο}

Λόγω της πολυπλοκότητας των σύγχρονων δικτύων, οι παραδοσιακές προσεγγίσεις ασφάλειας δεν επαρκούν πλέον. Απαιτούνται μηχανισμοί ανίχνευσης και απόκρισης σε πραγματικό χρόνο.

Τεχνολογίες όπως η μηχανική μάθηση και το \en{AI} χρησιμοποιούνται για την ανάλυση μεγάλων όγκων δεδομένων και τον εντοπισμό ανωμαλιών \cite{nguyen2018deep}.

Οι σύγχρονες λύσεις πρέπει να μπορούν να:
\begin{itemize}
    \item Ανιχνεύουν ανωμαλίες στην κίνηση δικτύου
    \item Αναγνωρίζουν κακόβουλα πρότυπα συμπεριφοράς
    \item Αντιδρούν αυτόματα σε απειλές
    \item Προσαρμόζονται σε εξελισσόμενες επιθέσεις
\end{itemize}

Στο \en{6G}, η έννοια της ασφάλειας βασισμένης σε \en{AI-native} προσεγγίσεις θα είναι κεντρική.

\subsection{Συμπέρασμα}

Η κυβερνοασφάλεια αποτελεί κρίσιμο παράγοντα για τα δίκτυα \en{5G} και \en{6G}, λόγω της πολυπλοκότητας και της σημασίας τους. Η εμφάνιση νέων απειλών και η διεύρυνση της επιφάνειας επίθεσης καθιστούν απαραίτητη την ανάπτυξη προηγμένων μηχανισμών ασφάλειας.

Με την αξιοποίηση τεχνολογιών όπως το \en{AI}, η ανάλυση σε πραγματικό χρόνο και τα αυτοματοποιημένα συστήματα απόκρισης στα δίκτυα επόμενης γενιάς μπορούν να επιτύχουν υψηλότερα επίπεδα ασφάλειας και ανθεκτικότητας. Το \en{6G} αναμένεται να διαδραματίσει καθοριστικό ρόλο στη διαμόρφωση ασφαλών και ευφυών υποδομών επικοινωνίας.